%======================================================================
\section{Conclusiones}\label{sec:conclusiones}
%======================================================================

\subsection{Respuestas a las preguntas}
\begin{enumerate}[label=\textbf{Q\arabic*.}]
  \item \textbf{¿Se puede resolver P0 con exactitud y en tiempo real?} Sí. La formulación
        monolítica resuelve P0 en $\sim0.05$~s (NLP) o $\sim10$~ms (QP condensado), unas
        $10^4$ veces más rápido que LJ y de forma reproducible.
  \item \textbf{¿Es el óptimo global?} La solución es un mínimo local estricto y el único
        óptimo hallado, y está demostradamente a menos del 2.3\,\% (1180~s) y del 11.4\,\%
        (1161~s) del óptimo global. La evidencia indica que la brecha real es $\approx0$.
  \item \textbf{¿Cuál método es mejor?} N3 y N2, empatados; N1 queda a 0.0026\,\%. LJ y GA
        cuestan al menos un 28\,\% y un 43\,\% más que el óptimo global.
  \item \textbf{¿El riesgo de pronóstico cambia la decisión?} No: su efecto sobre la PCT es
        unas 100 veces menor que la autoridad del actuador.
  \item \textbf{¿Las cotas del actuador son físicas?} No. Deben fijarse en kg/s, junto con una
        rampa admisible.
  \item \textbf{¿Se sostienen las conclusiones en lazo cerrado, con el HRM como planta?} Sí.
        Con la réplica del HRM en Julia (error $\sim10^{-6}$~K frente a Simulink), N1, N2 y N3
        dan el mismo resultado en lazo cerrado y en el Monte Carlo, sin fallas del solver. El
        desajuste ROM--HRM es un sesgo de $\sim0.2$~K que la corrección \emph{offset-free}
        elimina.
  \item \textbf{¿Resuelven a optimalidad?} N1 y N2 resuelven su QP al óptimo global, y N3
        llega a un óptimo local reproducible. Sus primeras acciones difieren en
        $\le0.11$~kg/s.
  \item \textbf{¿Dónde está el riesgo?} En los escenarios del modelo, en ninguna parte: la PCT
        no pasa de 1185~K. En la familia extrema, el resultado lo decide la capacidad del
        actuador, no el optimizador (sección~\ref{sec:extremos}).
\end{enumerate}

\begin{table}[H]\centering\small
\begin{tabularx}{\linewidth}{@{}l r r Y@{}}\toprule
Método & Tiempo por consulta & Costo en 1180~s & Garantía\\\midrule
LJ corregido (mejor semilla) & $\approx800$~s & 0.313 & ninguna; depende de la semilla\\
GA (mejor semilla) & $\approx615$~s & 0.349 & ninguna; depende de la semilla\\
N3, NLP monolítico & 0.05~s & 0.24403 & mínimo local estricto; a $\le2.3$\,\% del óptimo global\\
N2, QP iterado & 4 QP de 10~ms & 0.24403 & converge a un punto KKT de N3\\
N1, QP con una linealización & 0.01~s & 0.24404 & óptimo global de su QP (brecha dual $10^{-12}$)\\\bottomrule
\end{tabularx}
\caption{Resumen de los métodos en la consulta de 1180~s.}
\end{table}

\subsection{Lo que enseña este caso}
\begin{enumerate}
  \item \textbf{Conocer la estructura del problema vale más que la fuerza bruta.} Las
        metaheurísticas tratan el modelo como caja negra, y eso tiene sentido cuando el costo es
        discontinuo o ruidoso. En el horizonte del MPC, sin embargo, el costo es una función
        suave, y escribir el modelo como ecuaciones permite usar derivadas exactas.
  \item \textbf{La física explica la solución.} La función de transferencia muestra que el ACS
        actúa a través de un integrador. Eso explica por qué la decisión óptima es saturar al
        principio y luego seguir la referencia, y por qué el actuador tiene tan poca autoridad
        en 12.5~s.
  \item \textbf{Un modelo más simple puede bastar.} El QP con una sola linealización da
        prácticamente la misma decisión que el modelo no lineal completo, en una fracción del
        tiempo y con un certificado de optimalidad global para su propio problema.
  \item \textbf{Resolver bien un problema no es lo mismo que resolver el problema correcto.}
        La brecha de un solver dice si cumplió con su modelo. Para comparar decisiones hace
        falta una vara común, el costo no lineal, y una cota inferior demostrada.
  \item \textbf{Las garantías se construyen por capas}: condiciones locales, evidencia global
        y, por último, una cota inferior rigurosa con aritmética de intervalos.
  \item \textbf{Los supuestos heredados deben revisarse.} La cota del actuador en amperios
        parecía un dato, y resultó ser un artefacto del modelo de la bomba.
\end{enumerate}

\subsection{Limitaciones}
\begin{itemize}
  \item La certificación global se demostró en dos consultas. En lazo cerrado sólo hay
        evidencia de optimalidad local reproducible.
  \item El Monte Carlo tiene 6 realizaciones nominales/estresadas y 14 extremas. Bastan para
        las comparaciones pareadas con intervalos estrechos, no para estimar colas.
  \item La réplica del HRM se validó contra una sola realización archivada y en las salidas
        que usa el lazo. La validación nodo a nodo y con falla del ACS queda para leo\_asus.
  \item El pronóstico de fuga en el lazo es perfecto o tiene un sesgo constante; el ensamble
        BiLSTM todavía no está en el lazo.
  \item La certificación vale para P0 sobre el ROM. No dice nada sobre el desajuste entre el
        ROM y el HRM, ni sobre el reactor real.
  \item Las cotas del actuador siguen sin decidirse; se usan dos configuraciones de referencia.
  \item La brecha certificada es conservadora, sobre todo en 1161~s.
  \item La capa de riesgo todavía no se incorporó.
\end{itemize}

\subsection{Próximos pasos, por prioridad}
\begin{enumerate}
  \item \textbf{Fijar el actuador con el equipo}: caudal máximo y rampa admisible. En la
        familia extrema decide si se llega o no al estado final del DET.
  \item \textbf{Barrido de capacidad del actuador en la familia extrema}: buscar el régimen
        en que la capacidad es comparable a la fuga, donde la primera acción no está saturada
        y una falla del ACS puede llevar la PCT a 1220~K. Ahí la capa de riesgo puede cambiar
        decisiones.
  \item \textbf{Pronóstico en el lazo (F5b.0)}: modelo de error calibrado de la BiLSTM y
        ajuste paramétrico de la ley de Park. El ajuste fino de la BiLSTM, con datos de la
        réplica del HRM, sólo si las decisiones de riesgo resultan sensibles al pronóstico.
  \item \textbf{Capa de riesgo sobre N2 (F5b.1)}: escenarios de falla del ACS y de pronóstico
        (SAA con CVaR de la excedencia de la PCT) y restricción probabilística con margen de
        holgura analítico, comparadas en forma pareada contra N2+sesgo.
  \item \textbf{Campaña grande en Olab} ($\ge100$ realizaciones por familia).
  \item \textbf{Validación nodo a nodo del HRM} y corrida con falla del ACS en Simulink
        (leo\_asus).
\end{enumerate}
